\section{Documentación Detallada por Módulos}

\subsection{Módulo 1: Detección en Tiempo Real}

\subsubsection{Descripción General}

El módulo de detección proporciona visualización en tiempo real del procesamiento de video con YOLOv8, mostrando estadísticas actualizadas y permitiendo seleccionar entre múltiples cámaras.

\subsubsection{Componentes del Frontend}

\textbf{Archivo}: \texttt{resources/views/modulo1.blade.php} (323 líneas)

\paragraph{Estructura HTML Principal (Líneas 1-50)}

\begin{itemize}
    \item \textbf{Línea 1}: Extiende layout principal \texttt{@extends('layouts.app')}
    \item \textbf{Líneas 10-15}: Selector de cámara con ID \texttt{cameraSelector}
    \item \textbf{Línea 25}: Canvas para video \texttt{<canvas id="videoCanvas" width="800" height="600">}
    \item \textbf{Líneas 30-45}: Cards de estadísticas con IDs:
    \begin{itemize}
        \item \texttt{currentVehicles} - Vehículos actuales
        \item \texttt{totalVehicles} - Total detectados
        \item \texttt{fps} - Frames por segundo
        \item \texttt{averageVehicles} - Promedio de vehículos
    \end{itemize}
\end{itemize}

\paragraph{Lógica JavaScript (Líneas 150-323)}

\textbf{Variables Globales (Líneas 152-160)}:
\begin{lstlisting}[style=jsstyle, caption={Variables Globales - modulo1.blade.php}]
let websocket = null;
let currentCamera = 'CAM_002';
const canvas = document.getElementById('videoCanvas');
const ctx = canvas.getContext('2d');
let statsInterval = null;
let reconnectAttempts = 0;
const MAX_RECONNECT_ATTEMPTS = 5;
\end{lstlisting}

\textbf{Función connectWebSocket() (Líneas 170-210)}:
\begin{itemize}
    \item \textbf{Línea 171}: Construye URL WebSocket: \texttt{ws://localhost:8080/ws/stream?camera\_id=\$\{currentCamera\}}
    \item \textbf{Línea 173}: Crea instancia WebSocket con \texttt{new WebSocket(wsUrl)}
    \item \textbf{Línea 174}: Configura tipo binario: \texttt{websocket.binaryType = 'arraybuffer'}
    \item \textbf{Líneas 176-179}: Handler \texttt{onopen} - actualiza status a "Online"
    \item \textbf{Líneas 181-195}: Handler \texttt{onmessage}:
    \begin{itemize}
        \item Línea 182: Recibe ArrayBuffer del frame JPEG
        \item Línea 183: Convierte a Blob: \texttt{const blob = new Blob([event.data], \{type: 'image/jpeg'\})}
        \item Línea 184: Crea Object URL temporal
        \item Líneas 187-191: Dibuja imagen en canvas usando \texttt{ctx.drawImage()}
        \item Línea 192: Libera memoria con \texttt{URL.revokeObjectURL(url)}
    \end{itemize}
    \item \textbf{Líneas 197-204}: Manejo de errores y reconexión automática
\end{itemize}

\textbf{Función updateStats() (Líneas 215-250)}:
\begin{itemize}
    \item \textbf{Línea 217}: Fetch a API Laravel: \texttt{/api/vehicle-monitor/\$\{currentCamera\}}
    \item \textbf{Línea 218}: Parsea respuesta JSON
    \item \textbf{Líneas 220-223}: Actualiza elementos DOM:
    \begin{lstlisting}[style=jsstyle]
document.getElementById('currentVehicles').textContent = data.current_vehicles || 0;
document.getElementById('totalVehicles').textContent = data.total_detected || 0;
document.getElementById('fps').textContent = (data.fps || 0).toFixed(1);
    \end{lstlisting}
    \item \textbf{Líneas 225-230}: Llama \texttt{updateVehicleTypes(data.vehicle\_types)} para actualizar gráfico
    \item \textbf{Líneas 232-237}: Llama \texttt{updateVehicleColors(data.vehicle\_colors)} para badges de colores
\end{itemize}

\textbf{Inicialización (Líneas 300-315)}:
\begin{lstlisting}[style=jsstyle, caption={Inicialización del Módulo}]
window.addEventListener('DOMContentLoaded', () => {
    connectWebSocket();
    statsInterval = setInterval(updateStats, 1000); // Actualizar cada 1s
    
    document.getElementById('cameraSelector').addEventListener('change', (e) => {
        currentCamera = e.target.value;
        if (websocket) websocket.close();
        connectWebSocket();
    });
});
\end{lstlisting}

\subsubsection{Backend API}

\textbf{Archivo}: \texttt{app/Http/Controllers/VehicleMonitorController.php} (59 líneas)

\paragraph{Clase VehicleMonitorController}

\textbf{Propiedades (Líneas 12-13)}:
\begin{lstlisting}[style=phpstyle]
protected $detectorUrl = 'http://localhost:8080';
\end{lstlisting}

\textbf{Método getStats() (Líneas 20-57)}:
\begin{itemize}
    \item \textbf{Línea 21}: Recibe parámetro \texttt{\$cameraId} (default: 'CAM\_002')
    \item \textbf{Líneas 23-26}: Valida y normaliza cameraId:
    \begin{lstlisting}[style=phpstyle]
$cameraId = strtoupper($cameraId);
if (!in_array($cameraId, ['CAM_001', 'CAM_002'])) {
    $cameraId = 'CAM_002';
}
    \end{lstlisting}
    \item \textbf{Línea 28}: Construye URL completa: \texttt{"\{detectorUrl\}/api/vehicles?camera\_id=\{cameraId\}"}
    \item \textbf{Línea 29}: Hace petición HTTP con timeout de 5s: \texttt{Http::timeout(5)->get(\$url)}
    \item \textbf{Líneas 31-33}: Si exitoso, retorna JSON del detector Python
    \item \textbf{Líneas 35-42}: Si falla, retorna estructura por defecto con valores en 0
\end{itemize}

\subsubsection{Servicio Python Detector}

\textbf{Archivo}: \texttt{deteccion/vehiculo\_detector.py} (767 líneas)

\paragraph{API Flask - Endpoint /api/vehicles (Líneas 550-580)}

\textbf{Función get\_vehicle\_stats() (Línea 552)}:
\begin{lstlisting}[style=pythonstyle, caption={Endpoint de Estadísticas}]
@app.route('/api/vehicles', methods=['GET'])
def get_vehicle_stats():
    camera_id = request.args.get('camera_id', 'CAM_002')
    
    with data_lock:
        state = camera_states.get(camera_id, {})
        tracked = state.get('tracked_vehicles', {})
        
        # Contar por tipo de vehículo
        vehicle_types = {}
        vehicle_colors = {}
        
        for vehicle_id, vehicle_data in tracked.items():
            v_type = vehicle_data.get('type', 'desconocido')
            v_color = vehicle_data.get('color', 'desconocido')
            
            vehicle_types[v_type] = vehicle_types.get(v_type, 0) + 1
            vehicle_colors[v_color] = vehicle_colors.get(v_color, 0) + 1
        
        return jsonify({
            'camera_id': camera_id,
            'status': state.get('status', 'offline'),
            'current_vehicles': len(tracked),
            'total_detected': state.get('total_vehicles_detected', 0),
            'fps': state.get('fps', 0.0),
            'vehicle_types': vehicle_types,
            'vehicle_colors': vehicle_colors
        })
\end{lstlisting}

\paragraph{WebSocket Streaming - Endpoint /ws/stream (Líneas 600-650)}

\textbf{Función stream\_video() (Línea 602)}:
\begin{itemize}
    \item \textbf{Línea 604}: Obtiene \texttt{camera\_id} de query params
    \item \textbf{Líneas 608-615}: Loop infinito que:
    \begin{itemize}
        \item Línea 610: Accede a \texttt{camera\_states[camera\_id]['processed\_frame']} con lock
        \item Línea 612: Si frame existe, lo codifica a JPEG usando \texttt{cv2.imencode('.jpg', frame)}
        \item Línea 614: Envía bytes vía WebSocket: \texttt{ws.send(frame\_bytes)}
        \item Línea 616: Espera 33ms (30 FPS): \texttt{time.sleep(0.033)}
    \end{itemize}
    \item \textbf{Líneas 620-625}: Manejo de excepciones y cierre de conexión
\end{itemize}

\paragraph{Worker Thread - worker\_skyline() (Líneas 350-450)}

\textbf{Función principal (Línea 352)}:
\begin{itemize}
    \item \textbf{Líneas 355-365}: Llama \texttt{get\_skyline\_stream\_url\_robust()} para obtener URL HLS
    \item \textbf{Línea 370}: Abre stream con OpenCV: \texttt{cap = cv2.VideoCapture(stream\_url)}
    \item \textbf{Líneas 375-420}: Loop de procesamiento:
    \begin{itemize}
        \item Línea 377: Lee frame: \texttt{ret, frame = cap.read()}
        \item Línea 380: Llama \texttt{process\_frame\_generic(frame, 'CAM\_002')}
        \item Líneas 385-395: Calcula FPS usando \texttt{time.time()} y contador de frames
    \end{itemize}
\end{itemize}

\textbf{Función process\_frame\_generic() (Líneas 200-280)}:
\begin{itemize}
    \item \textbf{Línea 205}: Ejecuta YOLO: \texttt{results = model(frame, conf=0.4)}
    \item \textbf{Líneas 210-250}: Para cada detección:
    \begin{itemize}
        \item Línea 215: Obtiene bbox, confianza y clase
        \item Línea 220: Filtra solo vehículos (car, motorcycle, bus, truck)
        \item Línea 225: Llama \texttt{match\_vehicle\_to\_tracked(bbox, tracked, max\_distance=50)}
        \item Líneas 230-240: Si es nuevo vehículo:
        \begin{itemize}
            \item Asigna nuevo ID: \texttt{vehicle\_id = state['next\_vehicle\_id']}
            \item Incrementa contador
            \item Llama \texttt{get\_dominant\_color(frame, bbox)}
            \item Llama \texttt{save\_detection\_to\_xml(vehicle\_data, camera\_id)}
        \end{itemize}
        \item Línea 245: Dibuja bbox en frame con \texttt{cv2.rectangle()}
    \end{itemize}
    \item \textbf{Línea 270}: Almacena frame procesado: \texttt{state['processed\_frame'] = frame}
\end{itemize}

\subsubsection{Flujo Completo de Detección}

\begin{figure}[H]
\centering
\begin{tikzpicture}[
    box/.style={rectangle, draw, fill=cyan!15, minimum width=3.5cm, minimum height=0.7cm, font=\small},
    node distance=0.8cm
]
    \node[box] (step1) {1. Usuario abre modulo1.blade.php};
    \node[box, below=of step1] (step2) {2. DOMContentLoaded ejecuta (L.300)};
    \node[box, below=of step2] (step3) {3. connectWebSocket() (L.170)};
    \node[box, below=of step3] (step4) {4. WebSocket a localhost:8080/ws/stream};
    \node[box, below=of step4] (step5) {5. Python stream\_video() envía JPEGs};
    \node[box, below=of step5] (step6) {6. onmessage dibuja en canvas (L.181)};
    \node[box, below=of step6] (step7) {7. setInterval llama updateStats() (L.305)};
    \node[box, below=of step7] (step8) {8. Fetch a /api/vehicle-monitor};
    \node[box, below=of step8] (step9) {9. VehicleMonitorController::getStats()};
    \node[box, below=of step9] (step10) {10. HTTP a Python /api/vehicles};
    \node[box, below=of step10] (step11) {11. Retorna vehicle\_types, colors, FPS};
    \node[box, below=of step11] (step12) {12. Actualiza DOM (L.220-237)};
    
    \draw[-{Latex}] (step1) -- (step2);
    \draw[-{Latex}] (step2) -- (step3);
    \draw[-{Latex}] (step3) -- (step4);
    \draw[-{Latex}] (step4) -- (step5);
    \draw[-{Latex}] (step5) -- (step6);
    \draw[-{Latex}] (step6) -- (step7);
    \draw[-{Latex}] (step7) -- (step8);
    \draw[-{Latex}] (step8) -- (step9);
    \draw[-{Latex}] (step9) -- (step10);
    \draw[-{Latex}] (step10) -- (step11);
    \draw[-{Latex}] (step11) -- (step12);
\end{tikzpicture}
\caption{Flujo Completo del Módulo Detección}
\end{figure}

\newpage
\subsection{Módulo 2: Reportes y Estadísticas}

\subsubsection{Descripción General}

El módulo de reportes permite visualizar el historial completo de detecciones vehiculares con filtrados avanzados, gráficos estadísticos interactivos y exportación a Excel/CSV.

\subsubsection{Vista Principal - modulo2.blade.php}

\textbf{Archivo}: \texttt{resources/views/modulo2.blade.php} (341 líneas)

\paragraph{Estructura HTML (Líneas 1-180)}

\textbf{Filtros del formulario (L.24-65)}:
\begin{itemize}
    \item \textbf{L.27-30}: Select de cámara con opciones CAM\_001 y CAM\_002
    \item \textbf{L.35}: Input type='date' para fecha\_inicio, valor default: \texttt{date('Y-m-01')}
    \item \textbf{L.40}: Input type='date' para fecha\_fin, valor default: \texttt{date('Y-m-d')}
    \item \textbf{L.45-50}: Select de tipo con opciones: Todos, Auto, Moto, Bus
    \item \textbf{L.55-59}: Select de agrupación con modos: Automático, Por Hora, Por Día
\end{itemize}

\textbf{Canvas para gráficos (L.70-91)}:
\begin{itemize}
    \item \textbf{L.77}: Canvas \texttt{\#dailyChart} para gráfico de barras (300px alto)
    \item \textbf{L.89}: Canvas \texttt{\#typeDistributionChart} para gráfico de torta
\end{itemize}

\textbf{Tabla de registros (L.107-175)}:
\begin{lstlisting}[style=phpstyle, caption={Loop de Detecciones - modulo2.blade.php}]
@forelse($registros ?? [] as $registro)
    @php
        // L.120-126: Mapeo de colores por tipo
        $colorClass = match(strtolower($registro['tipo'])) {
            'auto' => 'bg-blue-500',
            'moto' => 'bg-red-500',
            'bus' => 'bg-green-500',
            default => 'bg-gray-500'
        };
        
        // L.127-128: Determinación de color según confianza
        $conf = floatval($registro['confianza']);
        $confColor = $conf > 90 ? 'bg-green-100' : 
                    ($conf > 75 ? 'bg-yellow-100' : 'bg-red-100');
        
        // L.131-142: Mapeo de colores de vehículo a CSS
        $vehicleColorMap = [
            'rojo' => 'bg-red-500',
            'azul' => 'bg-blue-500',
            'blanco' => 'bg-gray-100 border',
            // ... otros colores
        ];
    @endphp
    
    <tr class="hover:bg-gray-50">
        <td>{{ $registro['fecha'] }}</td>
        <td><span class="{{ $colorClass }}"></span> {{ $registro['tipo'] }}</td>
        <td><span class="{{ $colorBadge }}"></span> {{ $registro['color'] }}</td>
    </tr>
@empty
    <tr><td colspan="6">No se encontraron registros</td></tr>
@endforelse
\end{lstlisting}

\paragraph{JavaScript de Gráficos (L.186-340)}

\textbf{Variables globales (L.188-193)}:
\begin{lstlisting}[language=JavaScript]
let chartBar = null;    // Gráfico de barras
let chartPie = null;    // Gráfico de torta
const getSelectedCamera = () => {
    const params = new URLSearchParams(window.location.search);
    return params.get('camera_id') || 'CAM_002';
};
\end{lstlisting}

\textbf{Función processAndInitCharts() (L.209-333)}:

\textbf{Determinación de modo de visualización (L.211-231)}:
\begin{lstlisting}[language=JavaScript]
const viewMode = urlParams.get('view_mode') || 'auto';
let isMultiDay = false;

if (viewMode === 'daily') {
    isMultiDay = true;  // Agrupar por día
} else if (viewMode === 'hourly') {
    isMultiDay = false;  // Agrupar por hora 0-23
} else {
    // Modo Automático: detecta si hay múltiples días
    const uniqueDates = new Set(data.map(item => item.fecha.split(' ')[0]));
    if (uniqueDates.size > 1) isMultiDay = true;
}
\end{lstlisting}

\textbf{Procesamiento de datos (L.233-261)}:
\begin{lstlisting}[language=JavaScript]
const stats = {};
const typesCount = { 'Auto': 0, 'Moto': 0, 'Bus': 0, 'Otro': 0 };

data.forEach(item => {
    const [datePart, timePart] = item.fecha.split(' ');
    const hour = timePart.split(':')[0];
    
    // Clave de agrupación: fecha completa O hora
    const key = isMultiDay ? datePart : `${hour}:00`;
    
    // Inicializar contador si no existe
    if (!stats[key]) {
        stats[key] = { 'Auto': 0, 'Moto': 0, 'Bus': 0, 'Otro': 0 };
    }
    
    // Incrementar contador
    stats[key][item.tipo]++;
    typesCount[item.tipo]++;
});
\end{lstlisting}

\textbf{Creación de gráfico de barras (L.292-312)}:
\begin{lstlisting}[language=JavaScript]
chartBar = new Chart(ctxBar, {
    type: 'bar',
    data: {
        labels: sortedKeys,  // Ordenadas cronológicamente
        datasets: [
            { label: 'Autos', data: dataAuto, backgroundColor: '#3b82f6' },
            { label: 'Motos', data: dataMoto, backgroundColor: '#ef4444' },
            { label: 'Buses', data: dataBus, backgroundColor: '#10b981' }
        ]
    },
    options: {
        responsive: true,
        maintainAspectRatio: false,
        scales: {
            y: { beginAtZero: true, ticks: { stepSize: 1 } }
        }
    }
});
\end{lstlisting}

\subsubsection{Controlador ReporteController}

\textbf{Archivo}: \texttt{app/Http/Controllers/ReporteController.php} (171 líneas)

\paragraph{Método obtenerRegistrosFiltrados() (L.12-73)}

\textbf{Carga de XML (L.15-19)}:
\begin{lstlisting}[style=phpstyle]
$xmlPath = storage_path('app/vehiculos_db.xml');
if (file_exists($xmlPath)) {
    $xmlContent = file_get_contents($xmlPath);
    $xml = simplexml_load_string($xmlContent);
}
\end{lstlisting}

\textbf{Construcción de registro (L.23-30)}:
\begin{lstlisting}[style=phpstyle]
$registro = [
    'fecha' => (string)$det->fecha,
    'tipo'  => (string)$det->tipo,
    'confianza' => (float)$det->confianza,
    'color' => isset($det->color) ? (string)$det->color : 'desconocido',
    'camara' => isset($det->camara) ? (string)$det->camara : 'CAM_002',
    'nombre_camara' => isset($det->nombre_camara) ? 
        (string)$det->nombre_camara : 'Skyline Cochabamba'
];
\end{lstlisting}

\textbf{Aplicación de filtros (L.34-60)}:
\begin{itemize}
    \item \textbf{L.34-38}: Filtro por camera\_id - compara string exacto
    \item \textbf{L.40-44}: Filtro por tipo - compara string exacto
    \item \textbf{L.46-52}: Filtro por fecha inicio:
    \begin{lstlisting}[style=phpstyle]
$fechaRegistro = strtotime($registro['fecha']);
$fechaInicio = strtotime($request->fecha_inicio . ' 00:00:00');
if ($fechaRegistro < $fechaInicio) {
    $incluir = false;
}
    \end{lstlisting}
    \item \textbf{L.54-60}: Similar para fecha\_fin, añade '23:59:59'
\end{itemize}

\paragraph{Método exportarExcel() (L.82-140)}

\textbf{Setup de archivo (L.84-92)}:
\begin{lstlisting}[style=phpstyle]
$registros = $this->obtenerRegistrosFiltrados($request);
$nombreArchivo = 'reportes_' . date('Y-m-d_His') . '.csv';
$rutaTemporal = storage_path('app/temp');

if (!is_dir($rutaTemporal)) {
    mkdir($rutaTemporal, 0755, true);
}

$rutaCompleta = $rutaTemporal . '/' . $nombreArchivo;
\end{lstlisting}

\textbf{Escritura CSV con BOM UTF-8 (L.100-122)}:
\begin{lstlisting}[style=phpstyle]
$archivo = fopen($rutaCompleta, 'w');

// BOM para Excel UTF-8
fprintf($archivo, chr(0xEF).chr(0xBB).chr(0xBF));

// Encabezados
fputcsv($archivo, [
    'Fecha/Hora', 'Tipo Vehículo', 'Confianza (%)', 
    'Color', 'Cámara', 'Nombre Cámara'
], ',');

// Datos
foreach ($registros as $registro) {
    fputcsv($archivo, [
        $registro['fecha'],
        $registro['tipo'],
        number_format($registro['confianza'], 2),
        $registro['color'],
        $registro['camara'],
        $registro['nombre_camara']
    ], ',');
}

fclose($archivo);
\end{lstlisting}

\textbf{Descarga y auto-eliminación (L.127)}:
\begin{lstlisting}[style=phpstyle]
return response()->download($rutaCompleta, $nombreArchivo)
    ->deleteFileAfterSend(true);
\end{lstlisting}

\newpage
\subsection{Módulo 3: Automatizaciones y Administración}

\subsubsection{Descripción General}

Panel administrativo exclusivo para gestionar configuraciones de automatización, reglas de notificación, frecuencias de reportes y limpieza de datos XML.

\subsubsection{Controlador ModuloTresController}

\textbf{Archivo}: \texttt{app/Http/Controllers/ModuloTresController.php} (139 líneas)

\paragraph{Método index() (L.16-26)}

\textbf{Eager Loading con relaciones (L.18)}:
\begin{lstlisting}[style=phpstyle]
`$users = User::with('automationConfig')->get();
`\end{lstlisting}

Uso del método \texttt{with()} para cargar la relación \texttt{automationConfig} y evitar N+1 queries.

\textbf{Estadísticas agregadas (L.21-23)}:
\begin{lstlisting}[style=phpstyle]
`$totalNotifications = Notification::count();
`$totalReports = Report::count();
`$activeUsers = User::where('role', '!=', 'admin')->count();
\end{lstlisting}

\paragraph{Método saveConfiguration() (L.42-60)}

\textbf{Reglas de validación (L.44-53)}:
\begin{lstlisting}[style=phpstyle]
`$validated = $`request->validate([
    'traffic_threshold' => 'required|integer|min:1|max:1000',
    'traffic_minutes' => 'required|integer|min:1|max:60',
    'notify_high_traffic' => 'boolean',
    'notify_camera_offline' => 'boolean',
    'camera_offline_minutes' => 'required|integer|min:1|max:120',
    'report_frequency' => 'required|in:daily,weekly,monthly,disabled',
    'report_generation_time' => 'required|date_format:H:i',
    'auto_generate_reports' => 'boolean',
]);
\end{lstlisting}

\textbf{Actualización con Eloquent (L.55-56)}:
\begin{lstlisting}[style=phpstyle]
`$config = AutomationConfig::getOrCreateForUser(`$userId);
`$config->update(`$validated);
\end{lstlisting}

\subsection{Módulo 4: Accesibilidad (Comandos de Voz)}

\subsubsection{VoiceCommandController}

\textbf{Archivo}: \texttt{app/Http/Controllers/VoiceCommandController.php} (223 lined)

\paragraph{Método store() (L.32-60)}

Validación de 8 campos incluyendo pattern matching trigger.

\subsubsection{Servidor voice\_server.py}

\textbf{Archivo}: \texttt{deteccion/voice\_server.py} (180 líneas)

\paragraph{Clase VoiceRecognitionServer (L.20-145)}

\textbf{audio\_callback() (L.30-35)}: Captura audio a 16kHz, añade a cola.

\textbf{process\_audio() (L.37-85)}: Reconocimiento continuo con Vosk, emite \texttt{voice\_command} y \texttt{voice\_partial}.

\subsection{Módulo 5: Gestión de Usuarios}

\subsubsection{UserController}

Operaciones CRUD: crear con hash bcrypt, eliminar con validación de único admin, cambiar rol useradmin.

\section{Arquitectura Completa de la Aplicación}

\begin{figure}[H]
\centering
\begin{tikzpicture}[
    layer/.style={rectangle, draw, fill=blue!10, minimum width=11cm, minimum height=1.2cm, font=\small},
    component/.style={rectangle, draw, fill=green!15, minimum width=2.2cm, minimum height=0.7cm, font=\tiny},
    service/.style={rectangle, draw, fill=orange!15, minimum width=2cm, minimum height=0.7cm, font=\tiny},
    node distance=0.4cm and 0.3cm
]
    % Frontend Layer
    \node[layer] (frontend) {\textbf{CAPA DE PRESENTACIÓN - FRONTEND}};
    \node[component, below=of frontend, xshift=-4.5cm] (m1) {modulo1\\Detección};
    \node[component, right=of m1] (m2) {modulo2\\Reportes};
    \node[component, right=of m2] (m3) {modulo3\\Automatización};
    \node[component, right=of m3] (m4) {modulo4\\Voz};
    \node[component, right=of m4] (m5) {modulo5\\Gestión};
    
    % Backend Layer
    \node[layer, below=1.5cm of m3] (backend) {\textbf{CAPA DE APLICACIÓN - LARAVEL BACKEND}};
    \node[component, below=of backend, xshift=-4cm] (ctrl1) {VehicleMonitor\\Controller};
    \node[component, right=of ctrl1] (ctrl2) {Reporte\\Controller};
    \node[component, right=of ctrl2] (ctrl3) {ModuloTres\\Controller};
    \node[component, right=of ctrl3] (ctrl4) {VoiceCommand\\Controller};
    \node[component, right=of ctrl4] (ctrl5) {User\\Controller};
    
    % Services Layer
    \node[layer, below=1.5cm of ctrl3] (services) {\textbf{SERVICIOS Y LÓGICA DE NEGOCIO}};
    \node[service, below=of services, xshift=-2.5cm] (notif) {NotificationService};
    \node[service, right=of notif] (report) {ReportService};
    \node[service, right=of report] (auth) {AuthService};
    
    % Data Layer
    \node[layer, below=1.5cm of report] (data) {\textbf{CAPA DE DATOS}};
    \node[service, below=of data, xshift=-3.5cm] (mysql) {MySQL\\Database};
    \node[service, right=of mysql] (xml) {vehiculos\_db.xml};
    \node[service, right=of xml] (config) {users\_config.xml};
    
    % External Services
    \node[layer, below=1.5cm of xml] (external) {\textbf{SERVICIOS EXTERNOS - PYTHON}};
    \node[service, below=of external, xshift=-2cm] (detector) {vehiculo\_detector.py\\YOLOv8+Flask};
    \node[service, right=of detector] (voice) {voice\_server.py\\Vosk+SocketIO};
    
    % Connections
    \draw[-{Latex}] (m1) -- (ctrl1);
    \draw[-{Latex}] (m2) -- (ctrl2);
    \draw[-{Latex}] (m3) -- (ctrl3);
    \draw[-{Latex}] (m4) -- (ctrl4);
    \draw[-{Latex}] (m5) -- (ctrl5);
    
    \draw[-{Latex}] (ctrl2) -- (report);
    \draw[-{Latex}] (ctrl3) -- (notif);
    
    \draw[-{Latex}] (notif) -- (mysql);
    \draw[-{Latex}] (report) -- (mysql);
    \draw[-{Latex}] (report) -- (xml);
    \draw[-{Latex}] (ctrl2) -- (xml);
    \draw[-{Latex}] (ctrl3) -- (config);
    
    \draw[-{Latex}] (ctrl1) -- (detector);
    \draw[-{Latex}] (m1) -- (detector);
    \draw[-{Latex}] (detector) -- (xml);
    
    \draw[-{Latex}] (m4) -- (voice);
    
\end{tikzpicture}
\caption{Arquitectura Completa del Sistema PolySense}
\end{figure}

\end{document}
